% This document was generated by an LLM.

\documentclass{essay}

\DeclareMathOperator{\gd}{gd}

\title{Circular and Hyperbolic:\\
Trigonometric Functions, Their Hyperbolic Twins, and Their Inverses}
\author{}
\date{}

\begin{document}

\maketitle

\section{One function, two directions}

The circular functions ($\sin$, $\cos$, $\tan$) and the hyperbolic functions
($\sinh$, $\cosh$, $\tanh$) are closely related, but neither real family is a
special case of the other as a real-valued function. Extend their exponential
formulas to complex arguments and a quarter-turn of the argument converts one
family into the other. Circular sine oscillates on the real axis; on the
imaginary axis it grows like hyperbolic sine.

Many connections follow from this fact: conversion identities, logarithmic
formulas for the inverses, and paired integrals of calculus. There are also
real-variable connections through differential equations and area. One can
understand those first; complex analysis explains why the formulas continue to
hold after an imaginary substitution and how to handle their branches.

The essay starts with the exponential definitions, derives the conversion
identities, examines their geometry, inverts both families, and ends with the
consequences for calculus.

\section{The exponential definitions}

Euler's formula ties the circular functions to the exponential:
\[
  e^{iz} = \cos z + i \sin z .
\]
Write it at $z$ and at $-z$, then add or subtract the two equations. This gives
sine and cosine as combinations of $e^{iz}$ and $e^{-iz}$:
\[
  \cos z = \frac{e^{iz} + e^{-iz}}{2}, \qquad
  \sin z = \frac{e^{iz} - e^{-iz}}{2i}.
\]
The hyperbolic functions are the same combinations with the $i$ removed:
\[
  \cosh z = \frac{e^{z} + e^{-z}}{2}, \qquad
  \sinh z = \frac{e^{z} - e^{-z}}{2}.
\]
In both families $\cos$ and $\cosh$ are the even parts of an exponential, and
$\sin$ and $\sinh$ are the odd parts, up to a factor of $i$. Tangent, cotangent,
secant and the rest are quotients of these, so they inherit the same structure.

These formulas define all four functions on all of $\C$ as entire functions.
Their power series match that picture.
$\sinh z = z + z^3/3! + z^5/5! + \cdots$ has the same coefficients as $\sin z$,
without the alternating signs. The sign pattern in the sine series is exactly
what substituting $iz$ produces.

\section{A real-variable bridge: differential equations}

The two families can be compared before introducing complex numbers. On the
real line, their defining pairs satisfy
\[
  \frac{d}{dt}\binom{\cos t}{\sin t}
    = \begin{pmatrix}0&-1\\1&0\end{pmatrix}
      \binom{\cos t}{\sin t},
  \qquad
  \frac{d}{dt}\binom{\cosh t}{\sinh t}
    = \begin{pmatrix}0&1\\1&0\end{pmatrix}
      \binom{\cosh t}{\sinh t}.
\]
Both start at $(1,0)$. The sign in the upper-right corner distinguishes a
rotation from a hyperbolic rotation. In particular, each coordinate solves a
second-order equation:
\[
  y''=-y \quad\text{for }\sin,\cos,
  \qquad
  y''=y \quad\text{for }\sinh,\cosh.
\]
This is the oscillation-versus-growth distinction in its simplest form. The
corresponding addition laws are matrix multiplication laws:
\[
  \begin{pmatrix}\cos(s+t)&-\sin(s+t)\\
                  \sin(s+t)&\cos(s+t)\end{pmatrix}
  = R(s)R(t),
  \qquad
  \begin{pmatrix}\cosh(s+t)&\sinh(s+t)\\
                  \sinh(s+t)&\cosh(s+t)\end{pmatrix}
  = H(s)H(t),
\]
where $R$ and $H$ denote the displayed matrices at their respective arguments.
Both have determinant $1$; direct multiplication shows that they preserve
$x^2+y^2$ and $x^2-y^2$, respectively.
This gives the circle and hyperbola identities without complex analysis.

The imaginary substitution explains the changed sign: if $f''=-f$ and
$g(t)=f(it)$, the chain rule gives $g''(t)=-f''(it)=g(t)$. This short argument
uses the complex extension of $f$, but it reveals why $i^2=-1$ interchanges
the two differential equations.

\section{The rotation identities}

Substitute $iz$ into the circular definitions. Then $e^{i \cdot iz} = e^{-z}$,
and the circular formulas become the hyperbolic ones:
\[
  \cos(iz) = \cosh z, \qquad
  \sin(iz) = i \sinh z, \qquad
  \tan(iz) = i \tanh z .
\]
In the other direction:
\[
  \cosh(iz) = \cos z, \qquad
  \sinh(iz) = i \sin z, \qquad
  \tanh(iz) = i \tan z .
\]
The form $\sinh x = -i \sin(ix)$ is the first sine identity solved for $\sinh$.
The reciprocal functions follow by division: $\sec(iz) = \sech z$,
$\csc(iz) = -i \csch z$, $\cot(iz) = -i \coth z$.

The even functions ($\cos$, $\cosh$, $\sec$, $\sech$) carry over unchanged. The
odd ones pick up a factor of $i$. That is forced: multiplying the argument by
$i$ multiplies the $n$-th power-series term by $i^n$. For even $n$ this gives
$\pm 1$, which cancels against the alternating signs. For odd $n$ it leaves one
extra $i$.

Combined with the addition formulas, these identities split any circular
function into real and imaginary parts:
\begin{align*}
  \sin(x + iy) &= \sin x \cosh y + i \cos x \sinh y, \\
  \cos(x + iy) &= \cos x \cosh y - i \sin x \sinh y .
\end{align*}
One consequence is $|\sin(x + iy)|^2 = \sin^2 x + \sinh^2 y$. So sine is bounded
on the real line but grows like $e^{|y|}/2$ off it. Liouville's theorem requires
this, since a bounded entire function must be constant.

\paragraph{A reliable translation rule.}
For identities in sine and cosine, substitute $iz$ for each argument, replace
$\sin(iz)$ by $i\sinh z$ and $\cos(iz)$ by $\cosh z$, and simplify the powers of
$i$. Quotients can be handled the same way after writing out their numerators
and denominators. Some examples:
\begin{center}
\begin{tabular}{ll}
\toprule
Circular & Hyperbolic \\
\midrule
$\cos^2 x + \sin^2 x = 1$ & $\cosh^2 x - \sinh^2 x = 1$ \\
$\cos 2x = 1 - 2\sin^2 x$ & $\cosh 2x = 1 + 2\sinh^2 x$ \\
$1 + \tan^2 x = \sec^2 x$ & $1 - \tanh^2 x = \sech^2 x$ \\
$\sin(a+b) = \sin a \cos b + \cos a \sin b$
  & $\sinh(a+b) = \sinh a \cosh b + \cosh a \sinh b$ \\
\bottomrule
\end{tabular}
\end{center}
This is an identity substitution, not a mnemonic about signs: two sine factors
contribute $i^2=-1$, while four contribute $i^4=1$.

\section{Geometry: the circle and the hyperbola}

The point $(\cos t, \sin t)$ lies on the unit circle $x^2 + y^2 = 1$. The point
$(\cosh t, \sinh t)$ lies on the right branch of the unit hyperbola
$x^2 - y^2 = 1$. Substituting $y \mapsto iy$ turns one curve into the other,
which is the rotation identity seen in coordinates.

The common meaning of $t$ is signed area, not arc length. Sweep the ray from
the origin along the curve from parameter $0$ to $t$, counting orientation and
repeated turns. The swept area is $t/2$ in both cases. For a curve parameterized by
$(x(s), y(s))$, the area swept from the origin is
\[
  A = \frac{1}{2} \int_0^t \bigl( x(s)\,y'(s) - y(s)\,x'(s) \bigr)\, ds .
\]
For the circle the integrand is $\cos^2 + \sin^2 = 1$. For the hyperbola it is
$\cosh^2 - \sinh^2 = 1$. Both give $A = t/2$.

On the unit circle, the arc length is $|t|$, twice the unsigned sector area.
On the hyperbola, arc length is not $|t|$. That is why ``hyperbolic angle'' is
defined by twice the signed swept area. It is also why the
hyperbolic angle is unbounded, while circular angle wraps around after $2\pi$.

This shows the main structural difference between the two families. The circle
is compact, so the circular functions are periodic on $\R$. The hyperbola is
not, so the hyperbolic functions are injective or monotone on suitable pieces
of $\R$. In $\C$ the periods reappear rotated: $\sinh$ and $\cosh$ have period
$2\pi i$, and $\tanh$ has period $\pi i$.

\section{Why the inverses are logarithms}

Each of the six functions is a rational function of a single exponential:
$u = e^w$ for the hyperbolic family, $u = e^{iw}$ for the circular one. To solve
$f(w) = z$, rewrite it as a polynomial equation in $u$, solve for $u$, and take
a logarithm. For $\sin$, $\cos$, $\sinh$ and $\cosh$ the polynomial is
quadratic. For $\tan$ and $\tanh$ it is linear in $u^2$.

\paragraph{arcsinh.}
$\sinh w = z$ becomes $u - 1/u = 2z$, that is $u^2 - 2zu - 1 = 0$. The roots
multiply to $-1$, so for real $z$ exactly one is positive, namely
$z + \sqrt{z^2 + 1}$:
\[
  \arcsinh z = \ln\!\left( z + \sqrt{z^2 + 1} \right).
\]

\paragraph{arccosh.}
$\cosh w = z$ becomes $u^2 - 2zu + 1 = 0$. The roots multiply to $+1$, so they
are reciprocals and give $\pm w$. That is the evenness of $\cosh$. Taking the
larger root, for real $z \ge 1$:
\[
  \arccosh z = \ln\!\left( z + \sqrt{z^2 - 1} \right).
\]

\paragraph{arctanh.}
$\tanh w = (u^2 - 1)/(u^2 + 1) = z$, so $u^2 = (1+z)/(1-z)$:
\[
  \arctanh z = \frac{1}{2} \ln \frac{1 + z}{1 - z}.
\]

\paragraph{The circular inverses.}
These come out the same way with $u = e^{iw}$, and $w = -i \ln u$ at the end:
\begin{align*}
  \arcsin z &= -i \ln\!\left( iz + \sqrt{1 - z^2} \right), \\
  \arccos z &= -i \ln\!\left( z + i\sqrt{1 - z^2} \right), \\
  \arctan z &= \frac{1}{2i} \ln \frac{1 + iz}{1 - iz}.
\end{align*}

For real $z \in [-1, 1]$, the two roots of the arcsin quadratic are
$iz\pm\sqrt{1-z^2}$, and both lie on the unit circle. Choosing the
$+$ root puts $w$ in $[-\pi/2, \pi/2]$. The other root gives $\pi - w$, the
second real solution of $\sin w = z$. The periodicity $2\pi$ comes from the log,
which is defined only up to $2\pi i$; multiplied by $-i$, that becomes $2\pi$ in
$w$.

The pattern is general. Every inverse trigonometric and inverse hyperbolic
function is a logarithm of an algebraic function of $z$. Nothing new is needed
beyond $\log$ and the square root.

\section{Converting between the inverses}

Inverting the rotation identities gives conversion rules for the inverse
functions. Start from $\sin(iw) = i \sinh w$ and let $w = \arcsinh z$. Then
$\sin(iw) = iz$, so $iw = \arcsin(iz)$. The same argument run from the other
identity gives the second formula, and $\tan$ works the same way:
\begin{center}
\begin{tabular}{ll}
\toprule
Inverse rule & Comes from \\
\midrule
$\arcsin z = -i \arcsinh(iz)$ & $\sinh(iw) = i \sin w$ \\
$\arcsinh z = -i \arcsin(iz)$ & $\sin(iw) = i \sinh w$ \\
$\arctan z = -i \arctanh(iz)$ & $\tanh(iw) = i \tan w$ \\
$\arctanh z = -i \arctan(iz)$ & $\tan(iw) = i \tanh w$ \\
$\arccosh z = \pm i \arccos z$ & $\cosh(iw) = \cos w$ \\
\bottomrule
\end{tabular}
\end{center}
These hold for the principal branches. The cosine pair is the exception.
Because $\cos$ and $\cosh$ are even, inverting only fixes $iw$ up to sign. The
principal arccosh has non-negative real part, so the sign is $+i$ for
$\Im z > 0$ and $-i$ for $\Im z < 0$. On the real segment it depends on how
the cut is closed. For example, $\arccosh(1/2) = i\pi/3 = i \arccos(1/2)$.

The two routes to the logarithmic form now agree. Substituting $u = iz$ into
$\arcsinh u = \ln\bigl(u + \sqrt{u^2 + 1}\bigr)$ gives
\[
  \arcsin z = -i \arcsinh(iz) = -i \ln\!\left( iz + \sqrt{1 - z^2} \right).
\]
This is the same formula the direct derivation produced. So there is one
logarithmic formula per pair of functions, not per function. The rotation
identity transports it from one member of the pair to the other.

\section{Branches, cuts, and where the identities fail}

Every formula above contains two multivalued operations, a square root and a
logarithm. Choosing principal branches makes the inverses single-valued, at the
cost of a cut. The standard principal cuts are:
\begin{center}
\begin{tabular}{lll}
\toprule
Function & Cut & Values (strip) \\
\midrule
$\arcsin$, $\arccos$ & $(-\infty, -1] \cup [1, \infty)$
  & $\Re w \in [-\pi/2, \pi/2]$ and $[0, \pi]$ \\
$\arctan$ & $(-i\infty, -i] \cup [i, i\infty)$
  & $\Re w \in (-\pi/2, \pi/2)$ \\
$\arcsinh$ & $(-i\infty, -i] \cup [i, i\infty)$
  & $\Im w \in [-\pi/2, \pi/2]$ \\
$\arctanh$ & $(-\infty, -1] \cup [1, \infty)$
  & $\Im w \in (-\pi/2, \pi/2)$ \\
$\arccosh$ & $(-\infty, 1]$ & $\Re w \ge 0$ \\
\bottomrule
\end{tabular}
\end{center}
The rotation $z \mapsto iz$ turns the real cuts of arcsin into the imaginary
cuts of arcsinh. This is the inverse identity of the previous section, seen in
the cut geometry.

The identities fail in three typical ways.
\begin{enumerate}
  \item \textbf{Inverse after function is not the identity.}
    $\arcsin(\sin z) = z$ only in the principal strip. For example,
    $\arcsin(\sin \pi) = 0$.
  \item \textbf{Log rules break.}
    $\ln(ab) = \ln a + \ln b$ holds only up to $2\pi i$. So the arctanh formula
    cannot be split as $\tfrac{1}{2}\bigl(\ln(1+z) - \ln(1-z)\bigr)$ without
    checking the branches.
  \item \textbf{The square root moves the cut.}
    $\sqrt{z^2 - 1}$ and $\sqrt{z-1}\sqrt{z+1}$ agree when $\Re z > 0$ and
    differ by sign when $\Re z < 0$. The naive arccosh formula therefore gives
    the wrong branch near the negative axis. At $z=-2$, the naive argument
    $-2+\sqrt3$ has modulus $2-\sqrt3<1$. Any logarithm of it has negative
    real part $\ln(2-\sqrt3)$, whereas a principal $\arccosh$ boundary value
    has positive real part $\ln(2+\sqrt3)$. A formula consistent on the
    slit plane needs the factored root:
    \[
      \arccosh z = \ln\!\left( z + \sqrt{z + 1}\,\sqrt{z - 1} \right).
    \]
\end{enumerate}

The working rule: begin on a region where the branch choices are unambiguous,
then extend an identity only while both sides remain analytic on the same
connected region. Crossing a cut requires new branch information. Kahan's paper on branch cuts,
\emph{Much Ado About Nothing's Sign Bit}, is the standard reference for getting
these right in floating point.

\section{Calculus consequences}

The derivatives of the inverses fall into pairs that differ only in a sign
under the root or in the denominator:
\begin{center}
\begin{tabular}{lll}
\toprule
Function & Derivative & Real domain \\
\midrule
$\arcsin x$ & $1/\sqrt{1 - x^2}$ & $|x|<1$ \\
$\arcsinh x$ & $1/\sqrt{1 + x^2}$ & $x\in\R$ \\
$\arccosh x$ & $1/\sqrt{x^2 - 1}$ & $x>1$ \\
$\arctan x$ & $1/(1 + x^2)$ & $x\in\R$ \\
$\arctanh x$ & $1/(1 - x^2)$ & $|x|<1$ \\
\bottomrule
\end{tabular}
\end{center}
Read backwards, this explains why integrals that look alike produce such
different-looking answers. $\int dx/(1 - x^2)$ can be done by partial
fractions, which yields $\tfrac12\ln\lvert(1+x)/(1-x)\rvert$ on any interval
avoiding $\pm1$. On $(-1,1)$ this is $\arctanh x$. $\int dx/(1 + x^2)$ has no real
partial fractions, so the answer gets a new name, $\arctan$. Over $\C$ it
factors as $(1 + ix)(1 - ix)$, and partial fractions give exactly the
logarithmic formula for $\arctan$ from earlier. The ``new'' function is an
artifact of working over $\R$.

The same holds for the substitution families of integration:
\begin{center}
\begin{tabular}{lll}
\toprule
Integrand contains & Circular substitution & Hyperbolic substitution \\
\midrule
$\sqrt{a^2 - x^2}$ & $x = a \sin t$ & $x = a \tanh t$ \\
$\sqrt{a^2 + x^2}$ & $x = a \tan t$ & $x = a \sinh t$ \\
$\sqrt{x^2 - a^2}$ & $x = a \sec t$ & $x = a \cosh t$ \\
\bottomrule
\end{tabular}
\end{center}
The hyperbolic column often gives the answer directly in log form. The circular
column produces things like $\int \sec t \, dt = \ln\lvert \sec t + \tan t
\rvert$, which is a hyperbolic inverse in disguise.

\paragraph{The Gudermannian.}
A real-variable bridge between the families, with no $i$ involved, is the
Gudermannian function:
\[
  \gd x = \int_0^x \sech t \, dt = \arctan(\sinh x)
        = 2 \arctan\!\left( \tanh \frac{x}{2} \right).
\]
It satisfies $\sinh x = \tan(\gd x)$, $\cosh x = \sec(\gd x)$ and
$\tanh x = \sin(\gd x)$. Thus it links coordinates on the hyperbola to
trigonometric ratios at a real circular angle. For $|\theta|<\pi/2$ its inverse is
\[
  \gd^{-1}\theta=\int_0^\theta\sec u\,du
    =\arcsinh(\tan\theta).
\]
The integral also gives the vertical coordinate of the spherical Mercator
projection as a function of latitude $\theta$.

\end{document}
